\documentclass[11pt,a4paper]{article}
\usepackage[T1]{fontenc}
\usepackage[utf8]{inputenc}
\usepackage[margin=2.4cm]{geometry}
\usepackage{amsmath,amssymb,amsthm}
\usepackage{booktabs}
\usepackage{xcolor}
\usepackage{colortbl}
\usepackage{enumitem}
\usepackage[colorlinks,linkcolor=blue!50!black,citecolor=blue!50!black]{hyperref}
\usepackage[nameinlink,noabbrev]{cleveref}
\colorlet{cgood}{green!22}
\colorlet{cmid}{orange!30}
\colorlet{cbad}{red!22}
\newtheorem{proposition}{Proposition}
\newtheorem{lemma}[proposition]{Lemma}
\newtheorem{remark}[proposition]{Remark}
\newcommand{\meas}{\textsf{[M]}}
\newcommand{\infr}{\textsf{[I]}}
\newcommand{\code}[1]{\texttt{#1}}
\newcommand{\dt}{\Delta t}
\newcommand{\liq}{\ell}
\newcommand{\gas}{g}
\newcommand{\MM}{\mathcal M}
\newcommand{\W}{\mathsf W}
\newcommand{\E}{\mathsf E}
\newcommand{\Z}{\mathsf Z}
\newcommand{\R}{\mathsf R}
\newcommand{\C}{\mathsf C}
\newcommand{\Y}{\mathsf Y}
\newcommand{\G}{\mathsf G}
\newcommand{\ubar}{\bar u}
\newcommand{\g}{\cellcolor{cgood}}
\newcommand{\m}{\cellcolor{cmid}}
\newcommand{\rd}{\cellcolor{cbad}}
\title{Momentum transport in the Gram metric of the mimetic projection:\\
designs, identities and measurements}
\author{audit note, TwoPhaseAMR.jl}
\date{}
\begin{document}
\maketitle
\begin{abstract}
The mimetic projection (\code{MimeticProjection}, metric MMd) measures the phase velocities
in the Gram metric $\MM$; the momentum transport still lives in the lumped metric $\W$ of
the staggered phase masses.  This note asks for a transport that is consistent with $\MM$:
exact $\MM$-momentum, exact free-stream preservation with the volume-of-fluid (VOF) fluxes,
semi-discrete energy conservation, and no small-cell restriction of the time step.  The
answer has three parts.  (i)~These requirements are \emph{incompatible} for the Gram metric
as assembled: $\MM$ applied to a constant field gives the Gram tensors of the cells, whose
sum differs from the VOF volume by $O(h^2)$ and changes as the interface moves; exact
momentum and exact free stream together force that sum to be invariant
(\cref{prop:obstruction}), and a flux form with the VOF fluxes alone cannot preserve
constants at all (\cref{prop:flux}).  (ii)~Sliding the pressure points of the cut cells
along their symmetry lines by $O(h^3)$ (``mass matching'') makes the Gram mass equal to the
VOF mass and removes the obstruction.  (iii)~Within the remaining freedom, every design that
adds the change of the metric $\E=\MM-\W$ explicitly to the momentum (flux form, skew form,
skew redistribution) divides metric jumps by the inertia of slaved values and is unbounded
near small cut cells; the bounded design is the \emph{increment form}
$\MM^+\delta=\W^+\Delta u_\W$, which maps the momentum increment of the existing $\W$
transport through $\MM^{-1}$ like a force.  It preserves constants exactly in every sweep,
needs one solve with the band factorisation that the projection computes anyway, and its
increment is bounded by $\lambda_{\max}(\W,\MM)\le2.7$ \meas\ independently of the cut
volumes; a mass-proportional momentum restoration (DP) makes its $\MM$-momentum exact with
mass matching.  Exact semi-discrete energy is not reachable with a bounded design for the
present, piecewise-in-time, metric.  \Cref{sec:aperture} (added later) replaces the activity
rule ``staggered phase volume $>0$'' by ``phase aperture $>0$'': it removes the slaved values,
makes DG exact for constants to round-off and DP's restoration exact under the necessary mass
merge, and removes the energy excursion of a drop impact ($186E_0\to1.1E_0$); but births happen
at finite apertures, the metric stays discontinuous in time, the small ties remain necessary
for the projection, and DG, DS, DF remain unbounded.  \textbf{Open items}: R3 (energy) with a
bounded design; gas values on faces of aperture $10^{-5}$--$10^{-4}$ at density ratio 1000 stay
unbounded ($10^4$) with every design; the KH, RT, capillary-wave, combined-case and cost
measurements were not run.  \meas\ marks a measurement, \infr\ an inference.
Prototype: worktree \code{wt-mm-momentum}, branch \code{study/mm-momentum-transport};
runs and report in \code{mm-runs/mm\_momentum/}.
\end{abstract}

\tableofcontents

%======================================================================
\section{Setting and requirements}\label{sec:setting}
\emph{Unknowns.} The phase values $u\in\mathbb R^N$ are the degrees of freedom of the
assembly of the projection: one value per face and phase with positive staggered phase
volume, one identified value on a grid-line interface face.  At a time level the geometry
(PLIC reconstruction, shared apertures, pressure points) gives
\begin{itemize}[leftmargin=*,itemsep=1pt]
\item the Gram metric $\MM$: per cut cell and phase the block
  $\rho^\pi h^2(\R\G^{+}\R^\top+\mathsf S)$ on the wetted sides, $\tfrac12\rho^\pi h^2$ per
  side of an uncut cell, ties of slaved and small-aperture values to donors, lumped masses
  where no block gives inertia (\code{\_assemble\_cut\_cell\_metric!});
\item the lumped metric $\W=\operatorname{diag}(m)$ with the staggered phase masses
  $m_{(f,\pi)}=\rho^\pi|V_f|\,\bar\alpha^\pi_f$, $\bar\alpha^\pi$ the interpolated VOF
  fractions; the code calls them \code{dof\_mass};
\item the excess $\E=\MM-\W$, which vanishes outside the band of cut cells (an uncut face
  between two pure cells of one phase has $\MM_{ff}=\rho h^2=m_f$);
\item the constant fields $\Z=[\Z^\liq_x,\Z^\liq_y,\Z^\gas_x,\Z^\gas_y]$, $(\Z^\pi_k)_{(f,\pi')}=1$
  if $f$ is normal to $e_k$ and $\pi'\in\{\pi,\text{identified}\}$, else $0$.
\end{itemize}
The $\MM$-momentum is $P=\Z^\top\MM u$ (per phase and component; the mixture momentum sums
the phases), the $\MM$-energy $\tfrac12u^\top\MM u$.  The projection conserves the mixture
$P$ up to the wall force and does not increase the $\MM$-energy.

\emph{The current transport} (techref ch.~6, \S6.8) advances each phase value in two
directional Lax--Wendroff sweeps with the dual phase volume fluxes interpolated from the VOF
fluxes, in fluctuation form with the predicted staggered fractions.  For this note it is the
map $u\mapsto u_\W$ with the properties used below \infr\ (techref eqs.\ momentum:conservative,
momentum:mass-identity, momentum:skew):
\begin{equation}\label{eq:w-transport}
  \W^+u_\W=\W u-\dt\,\C_\W(u),\qquad \Z^\top\C_\W(u)=0,\qquad
  \dt\,\C_\W(\Z)=-(\W^+-\W)\Z ,
\end{equation}
the last two per sweep (conservation and the mass identity), and for the all-centred
states in the semi-discrete limit $\C_\W+\C_\W^\top=-\dot\W$.  The CTU switch, the drain
limiter and the fraction floor modify the states near small fractions; constants are
preserved by all of them (fluctuation form).

\emph{Requirements} (the owner's, in the order of the assignment).
\begin{description}[leftmargin=1.1cm,style=nextline,itemsep=2pt]
\item[R1] $\MM$-momentum: $\Z^\top\MM^+u^+=\Z^\top\MM u$ on periodic or no-flux domains.
\item[R2] Free stream with the VOF fluxes: the momentum fluxes are built from the VOF
  volume fluxes (mapped to the dual control volumes as now); a constant field stays
  constant to round-off through every sweep; the momentum of a constant field is
  $\rho^\pi\times$ VOF volume $\times\,u$.
\item[R3] Energy: $\frac{\mathrm d}{\mathrm dt}\tfrac12u^\top\MM u=0$ for the centred states,
  so that Lax--Wendroff dissipates $O(\dt)$ only.
\item[R4] No small-cell restriction: the full-cell CFL bound of the current transport
  suffices for any cut volume; bounded phase values near tiny cut volumes.
\end{description}

%======================================================================
\section{What the Gram metric does to a constant field}\label{sec:constant}
\begin{proposition}[Constant fields]\label{prop:constant}
Let every cut cell have a consistent block, $m\mathsf N=\R$ (Prop.~2 of the note
\code{metric\_options}), and let no value be lumped.  Then
\begin{enumerate}[leftmargin=*,itemsep=1pt,label=(\alph*)]
\item $(\MM\Z^\pi_k)_{(f,\pi)}=\rho^\pi a^\pi_f h\,(x^*_{c_+(f)}-x^*_{c_-(f)})_k$, for faces of
  both orientations (``targets'' $\R$; on a face normal to $e_t$, $t\ne k$, the entry is the
  tangential offset of the two pressure points);
\item $\Z^{\pi\top}\MM\Z^\pi=\rho^\pi\sum_c\G^\pi_c$ ($2\times2$ per phase), with $\G^\pi_c=h^2\mathsf I$
  or $0$ in an uncut cell and the Gram tensor at the pressure point in a cut cell, if no
  value is identified (an identified grid-line value adds the diagonal entry of the other
  phase to both columns);
\item $\Z^{\pi\top}\W\Z^\pi=\rho^\pi V^\pi\,\mathsf I$ on a uniform periodic mesh ($|V_f|=h^2$,
  face fractions the means of the two cell fractions), $V^\pi$ the VOF volume of the phase.
\end{enumerate}
\end{proposition}
\begin{proof}
(a) is \code{metric\_options} Prop.~3; ties vanish on constants because a value and its
donor are parallel faces (the face ring).  (b) Sum (a) against $\Z^\pi$: the side terms of
a cell sum to its Gram tensor (\code{metric\_options} Prop.~4).  (c) The interpolated
fractions of the faces normal to $e_k$ sum to the cell fractions, each cell contributing
half of its volume to each of its two faces.
\end{proof}
So the $\MM$-mass of a phase is $\rho^\pi\sum_c\G^\pi_c$, not $\rho^\pi V^\pi$.  For a
closed aperture polygon with the pressure point at the chord midpoint
$\G^\liq_c=|P_c|\,\mathsf I$, the polygon area; with apertures averaged from two traces
$|P_c|\neq\alpha_ch^2$, and the pressure point is moved within its admissible interval.
Per value the two metrics differ by $O(1)$: a slaved value (aperture $0$) has
$(\MM\Z)_f=0$ and $(\W\Z)_f=m_f$, and the cross entries of (a) have no counterpart in $\W$.
\meas\ (\code{geometry.jl}, drop $R=0.2$, $\rho=(1,10^{-3})$, sliding through one cell):
\begin{center}\footnotesize
\setlength\tabcolsep{4pt}
\begin{tabular}{@{}rccccc@{}}
\toprule
$n$ & mean $\frac{(\sum_c\G^\liq_c)_{xx}-V^\liq}{V^\liq}$ & range (slide) &
$\frac{|(\sum_c\G^\liq_c)_{xy}|}{V^\liq}$ & mass-matched & max $|s_c|$\\
\midrule
16 & $-3.5\times10^{-3}$ & $6.4\times10^{-3}$ & $7.1\times10^{-4}$ & $6.6\times10^{-16}$ & $3.1\times10^{-2}h$\\
32 & $-4.2\times10^{-4}$ & $6.9\times10^{-4}$ & $3.6\times10^{-5}$ & $8.8\times10^{-16}$ & $7.7\times10^{-3}h$\\
64 & $-1.1\times10^{-4}$ & $8.3\times10^{-5}$ & $3.3\times10^{-6}$ & $6.6\times10^{-16}$ & $2.5\times10^{-3}h$\\
\bottomrule
\end{tabular}
\end{center}

The mismatch is $O(h^2)$ relative and moves by $O(h^3)$ as the drop slides.

\emph{Mass matching.} Moving $x^*$ along the symmetry line by $s\hat\eta$ changes
$\G^\liq$ by $-s|\eta|\hat\eta\hat\eta^\top$ and $\G^\gas$ by the opposite
(\code{metric\_options} Prop.~2(c)).  The slides of least $\sum_cs_c^2$ with
$\sum_c\G^\liq_c(s_c)=\sum_c\alpha_ch^2\,\mathsf I$ are
\begin{equation}\label{eq:match}
  s_c=|\eta_c|\,\langle\hat\eta_c\hat\eta_c^\top,\Lambda\rangle,\qquad
  \sum_c|\eta_c|^2\langle\hat\eta_c\hat\eta_c^\top,\Lambda\rangle\,\hat\eta_c\hat\eta_c^\top
  =\sum_c\bigl(\G^\liq_c(0)-\alpha_ch^2\mathsf I\bigr),
\end{equation}
a $3\times3$ symmetric system for $\Lambda$.  Its rank is the number of distinct normal
directions $\pm\hat\eta_c$ (at most three): with two or fewer (a straight interface, or normals
only along the axes) only the part of the right side in $\operatorname{span}\{\hat\eta_c\hat\eta_c^\top\}$
can be matched, and the remainder stays (it vanishes for a straight interface with exact
apertures, not in general for averaged traces).  When it is solvable, $\Z^\top\MM\Z=\Z^\top\W\Z$
exactly (both phases, since
$\G^\liq+\G^\gas=h^2\mathsf I$ and $\alpha^\liq+\alpha^\gas=1$).  The slides are $O(h^3)$
\meas\ (table above, ``max slide'').  This is the statement R2 asks for, in total: \emph{with
mass matching, $\MM$ applied to a constant field carries $\rho^\pi$ times the VOF volume of
the phase}; per value it cannot hold (\cref{prop:flux}).  It costs one global $3\times3$
solve per assembly (serial sums: bitwise reproducible) and touches the projection through
the pressure points (accuracy in \cref{sec:results}).

%======================================================================
\section{Obstructions}\label{sec:obstruction}
Consider any update of the form
\begin{equation}\label{eq:general}
  \MM^+u^+=\MM u-L(u),\qquad L\ \text{linear}.
\end{equation}
\begin{proposition}[Momentum and free stream need an invariant Gram mass]\label{prop:obstruction}
\eqref{eq:general} conserves $\Z^\top\MM u$ for every $u$ iff $\Z^\top L=0$, and preserves
every constant field $u=\Z U$ iff $L\Z=-(\MM^+-\MM)\Z$.  Both together imply
$\Z^\top(\MM^+-\MM)\Z=0$, i.e.\ $\sum_c\G^\pi_c$ is the same at both time levels.  With the
VOF mass conserved this is $\Z^\top(\E^+-\E)\Z=0$.
\end{proposition}
\begin{proof}
The first two statements are \eqref{eq:general} multiplied by $\Z^\top$ and evaluated at
$\Z U$.  Then $0=\Z^\top L\Z=-\Z^\top(\MM^+-\MM)\Z$; \cref{prop:constant}(b),(c).
\end{proof}
By \cref{sec:constant} the obstruction is nonzero for the Gram metric as assembled
\meas: no transport can satisfy R1 and R2 exactly.  With mass matching it vanishes.  The
condition is sharp: if $\Z^\top(\MM^+-\MM)\Z=0$, the dense $L=DP^\top$ with $D=-\Delta\MM\Z$,
$P=\Z(\Z^\top\Z)^{-1}$ satisfies both (locality is the issue, \cref{prop:redistribution}).

\begin{proposition}[A flux form with the VOF fluxes cannot preserve constants]\label{prop:flux}
Let $L=\dt\,\C_\W$, the momentum fluxes of the VOF volume fluxes \eqref{eq:w-transport}.
Then \eqref{eq:general} preserves constants iff $(\MM^+-\W^+)\Z=(\MM-\W)\Z$ per value, i.e.\
$\E\Z$ is invariant.  A value that enters or leaves the band (where $\E\Z=0$) must then have
$(\E\Z)_f=0$ already; for a moving interface every band value does so in turn, which
requires $\MM\Z=\W\Z$ on the band and fails by $O(1)$ per value (\cref{sec:constant}).  (For
a static geometry the condition is trivially met.)
\end{proposition}
\begin{proof}
$\MM^+\Z U=\MM\Z U-\dt\C_\W(\Z U)=\MM\Z U+(\W^+-\W)\Z U$ by \eqref{eq:w-transport}.
\end{proof}
Hence any $\MM$-consistent design with exact free stream must add terms built from the
change of the metric, $\E^+-\E$, to the VOF fluxes; and the fluxes of R2 transport $\W$-mass,
so the question is only how the band term $\E$ enters.

%======================================================================
\section{Designs}\label{sec:designs}
All designs keep the existing $\W$ sweeps (VOF fluxes, Lax--Wendroff, CTU switch, drain
limiter, floor) unchanged and add one solve with $\MM^+$ at the end of the step; $\ubar$
denotes the extended start values on the new degrees of freedom (the frozen sources of the
sweeps), $u_\W$ the values after the sweeps, and $c=\E u$ the old excess covector mapped
to the new values (an identified value split by its mass weight; values that disappear are
``orphans'').
\begin{description}[leftmargin=1.1cm,style=nextline,itemsep=2pt]
\item[D0] current: $u^+=u_\W$.
\item[D$\theta$] $\MM^+u^+=\MM^+\ubar+\W^+(u_\W-\ubar)+(1-\theta)\,(c-\E^+\ubar)$;
  \textbf{DI} $\theta=1$ (increment form), \textbf{DS} $\theta=\tfrac12$ (skew form),
  \textbf{DF} $\theta=0$ (flux form: $\MM^+u^+=\MM u-\dt\C_\W(u)$).
\item[DG] DS plus a skew band coupling $\Y$, $\Y\Z=\tfrac12(\E^+\Z-(\E\Z)_{\text{mapped}})$
  (least squares, minimal norm, on the graph of same-phase pairs of a cell, of the metric
  and of the face ring), plus the remainder $\Psi(\ubar)=\sum_k r_k\langle\ubar\rangle_k$ that
  keeps constants exact ($r=\tfrac12\Delta\E\Z-\Y\Z$, $\langle\cdot\rangle_k$ the mass mean of
  column $k$ over the ring).
\item[DP] DI plus a mass-proportional momentum restoration: with the update written as
  $\MM^+u^+=\MM^+\ubar+r$, the $\MM$-momentum change less its constant part,
  $d_k=\Z_k^\top(\MM^+\ubar+r)-\Z_k^\top\MM u-\sum_j\Z_k^\top\Delta\MM\Z_j\langle\ubar\rangle_j$
  ($=\Z_k^\top(\E^+\ubar-\E u)-\sum_j\Z_k^\top\Delta\E\Z_j\langle\ubar\rangle_j$ for DI with conservative
  sweeps), is removed by the force $-\W\Z_k\,d_k/\Z_k^\top\W\Z_k$ (identified values excluded).
  DGP is DG with the same restoration.
\end{description}
Each needs the excess of the old level (one product with the cached $\MM$ before the
sweeps) and one solve with the band factorisation of $\MM^+$ (\code{\_mimetic\_metric\_solve!}),
which the viscous step and the projection of the same step reuse; DG adds an LSQR solve on
the band graph.

\begin{proposition}[The $\theta$ family]\label{prop:theta}
Let \eqref{eq:w-transport} hold and let no value disappear.  Then for D$\theta$
\begin{enumerate}[leftmargin=*,itemsep=1pt,label=(\alph*)]
\item constants: $\MM^+(u^+-\Z U)=-(1-\theta)(\E^+-\E)\Z U$: exact iff $\theta=1$;
\item momentum: $\Z^\top\MM^+u^+-\Z^\top\MM u=\theta\,\Z^\top(\E^+-\E)u$: exact iff
  $\theta=0$;
\item energy, semi-discrete with centred states:
  $\frac{\mathrm d}{\mathrm dt}\tfrac12u^\top\MM u=(\theta-\tfrac12)\,u^\top\dot\E u$:
  exact iff $\theta=\tfrac12$.
\end{enumerate}
The three properties are governed by three different defects, $\Delta\E\Z$ (a), $\Z^\top\Delta\E u$
(b) and $u^\top\dot\E u$ (c); when all three are nonzero (generic in the band, measured) no
$\theta$ satisfies two of R1--R3.  When the degrees of freedom change, (b) reads exactly
$\theta(\Z^{+\top}\E^+\ubar-\Z^\top\E u)-(1-\theta)\,\mathrm{orph}$ with the excess of the
values that disappear.
\end{proposition}
\begin{proof}
Write the update as $\MM^+u^+-\MM u=\W^+u_\W-\W u+\theta(\E^+-\E)u$ (expand
$\MM^+\ubar=\W^+\ubar+\E^+\ubar$).  (a) For $u=\Z U$, $u_\W=\Z U$, so the right side is
$(\W^+-\W)\Z U+\theta\Delta\E\Z U=\Delta\MM\Z U-(1-\theta)\Delta\E\Z U$.  (b) Multiply by
$\Z^\top$ and use $\Z^\top\W^+u_\W=\Z^\top\W u$.  (c) In the limit,
$\frac{\mathrm d}{\mathrm dt}(\MM u)=-\C_\W u+\theta\dot\E u$, so
$\frac{\mathrm d}{\mathrm dt}\tfrac12u^\top\MM u=u^\top\frac{\mathrm d}{\mathrm dt}(\MM u)-\tfrac12u^\top\dot\MM u
=-u^\top\C_\W u+\theta u^\top\dot\E u-\tfrac12u^\top\dot\W u-\tfrac12u^\top\dot\E u$ and
$u^\top\C_\W u=-\tfrac12u^\top\dot\W u$.
\end{proof}
\meas\ (\code{smoke.jl}): one step of each design on a drop, $\rho=(1,10^{-3})$, $n=32$,
reproduces (b) to $2\times10^{-17}$ absolute (momentum $\approx0.13$), orphans included.

\begin{proposition}[Skew redistribution]\label{prop:redistribution}
If $\Z^\top(\E^+-\E)\Z=0$, there is a skew $\Y$ with $\Y\Z=\tfrac12(\E^+-\E)\Z$, and DG
with $\Psi=0$ satisfies (a), (b) and (c) of \cref{prop:theta} exactly.  On a graph whose
edges join the values of a phase in the cut cells, the kernel of the edge-to-value map
$\Y\mapsto\Y\Z$ (transposed) contains the symmetric $2\times2$ fields
$\{\Z\mathsf S\}$ per connected component (it equals them only if the $x$- and the $y$-values
of a component are each connected through same-axis edges), so a local $\Y$ needs the
condition per component: one global mass match does not suffice when the band has several
components (checked counterexample in \code{proof\_review/mimetic-momentum-transport}).
\end{proposition}
\begin{proof}
Existence: with $D=\tfrac12\Delta\E\Z$ and $P=\Z(\Z^\top\Z)^{-1}$, $\Y=DP^\top-PD^\top$ is
skew and $\Y\Z=D-P(\Z^\top D)^\top=D$ because $\Z^\top D=0$.  (a)
$-\tfrac12\Delta\E\Z+\Y\Z=0$.  (b) $\Z^\top(\tfrac12\Delta\E u+\Y u)=\tfrac12(\Delta\E\Z)^\top u-(\Y\Z)^\top u=0$
($\Delta\E$ symmetric, $\Y$ skew).  (c) $u^\top\Y u=0$ and the $\dot\E$ terms of
\cref{prop:theta}(c) cancel at $\theta=\tfrac12$.  Kernel: an edge $(f,g)$ maps $y$ to $+y$
in row $(f,\mathrm{axis}(g))$ and $-y$ in row $(g,\mathrm{axis}(f))$; $r_{(f,k)}=\mathsf S_{\mathrm{axis}(f),k}$
with $\mathsf S$ symmetric is annihilated.
\end{proof}

\begin{proposition}[Increment form: bounded, no small-cell restriction]\label{prop:increment}
The DI increment $\delta=u^+-\ubar=(\MM^+)^{-1}\W^+(u_\W-\ubar)$ satisfies
\[
  \|\delta\|_{\MM^+}^2\le\lambda_{\max}(\W^+,\MM^+)\,\|u_\W-\ubar\|_{\W^+}^2,\qquad
  \|\delta\|_\infty\le\|(\MM^+)^{-1}\W^+\|_\infty\,\|u_\W-\ubar\|_\infty ,
\]
and preserves constants exactly in every sweep.  The time step enters only through the
$\W$ sweeps, which are bounded under the full-cell CFL of the current transport.
\end{proposition}
\begin{proof}
$\|\delta\|_\MM^2=v^\top\W\MM^{-1}\W v\le\lambda_{\max}(\MM^{-1/2}\W\MM^{-1/2})\,v^\top\W v$
with $v=u_\W-\ubar$.  Constants: $u_\W=\ubar=\Z U$.
\end{proof}
\meas\ over random drops and lines with cut volumes down to $10^{-12}$ (\code{geometry.jl}):
\begin{center}\footnotesize
\setlength\tabcolsep{4pt}
\begin{tabular}{@{}lrlccccc@{}}
\toprule
geometry & $n$ & floor & min $\min(\alpha,1-\alpha)$ & $\lambda_{\max}(\W,\MM)$ &
$\|\MM^{-1}\W\|_\infty$ & $\min(\MM^{-1}\W)$ & $\max_f\mu_f$\\
\midrule
drop & 16 & $\sqrt{\epsilon}$ & $9.5\times10^{-8}$ & 2.63 & 4.98 & $-0.80$ & 1.00\\
drop & 32 & $\sqrt{\epsilon}$ & $5.3\times10^{-8}$ & 2.54 & 4.97 & $-0.81$ & 1.00\\
drop & 16 & $10^{-12}$ & $3.0\times10^{-12}$ & 2.63 & 4.95 & $-0.77$ & 1.00\\
drop & 32 & $10^{-12}$ & $4.3\times10^{-12}$ & 2.41 & 4.53 & $-0.68$ & 1.00\\
lines & 16 & $\sqrt{\epsilon}$ & $1.0\times10^{-6}$ & 2.22 & 4.28 & $-0.25$ & 1.00\\
lines & 32 & $\sqrt{\epsilon}$ & $1.7\times10^{-8}$ & 2.18 & 4.00 & $-0.25$ & 1.00\\
lines & 16 & $10^{-12}$ & $1.3\times10^{-12}$ & 2.22 & 4.28 & $-0.25$ & 1.00\\
lines & 32 & $10^{-12}$ & $1.0\times10^{-12}$ & 2.06 & 3.77 & $-0.25$ & 1.00\\
\bottomrule
\end{tabular}
\end{center}
40 random configurations per row (radius, sliver depth $10^{-1}$--$10^{-9}h$, line
direction), worst value over the row, $\rho=(1,10^{-3})$;
$\mu_f=|((\MM-\W)\Z)_f|/\max(m_f,\MM_{ff})$ is the per-value mismatch of the constant fields.

The algebra alone gives no cut-volume independence ($\lambda_{\max}\ge\max_f m_f/\MM_{ff}$); it
comes from the assembly: the tie of a slaved value carries its own staggered mass, so the row
of $\MM$ scales with the row of $\W$ (measured, no analytic bound).  The entries of
$\MM^{-1}\W$ are not all non-negative (min $\approx-0.8$): the map is bounded but not
monotone, and DI can create small new extrema near the interface.

\begin{remark}[Why the designs with $\E$ terms are unbounded]\label{rem:unbounded}
DF, DS and DG add $(\MM^+)^{-1}(c-\E^+\ubar)$ (and $\Y\ubar$).  The row of a slaved value
(aperture $0$) of $\MM^+$ is its tie, of weight equal to its staggered mass, as small as
$10^{-14}$ \meas; the same row of $c-\E^+\ubar$ carries the excess of the old level, which
belonged to a value with a Gram block or another donor.  $\E$ changes by $O(1)$, not
$O(\dt)$, when a value changes its status (wetted~$\leftrightarrow$~slaved, donor switch,
lumping, grid-line identification, cut~$\leftrightarrow$~uncut), so the quotient is not
small for small $\dt$.  For constants these terms cancel exactly
(\cref{prop:theta}(a), \cref{prop:redistribution}); for any other field they do not.
\meas: one to five steps with a vortex perturbation gave corrections $|u^+-u_\W|$ of $64$
(DS), $1.2\times10^{3}$ (DF), $4\times10^{3}$--$6\times10^4$ (DG), all on slaved values,
against $0.13$ for DI (\code{smoke.jl}).
\end{remark}

\begin{proposition}[Momentum of the increment form]\label{prop:di-momentum}
For DI, $P^{N}-P^{0}=\Z^\top\E^{N}u^{N-1}-\Z^\top\E^{0}u^{0}-\sum_{n=1}^{N-1}\Z^\top\E^{n}(u^{n}-u^{n-1})$
(plus the $\W$-momentum defects of the sweeps): the drift is bounded by $\|\E\|$ times the
variation of $u$ in time and does not grow for a steady field.  With mass matching the
constant part vanishes, and DP removes the rest: $\Z^\top\MM^+u^+=\Z^\top\MM u$ exactly,
constants exact.
\end{proposition}
\begin{proof}
Sum \cref{prop:theta}(b) at $\theta=1$ by parts.  DP: the restoration force has
$\Z_k^\top$-moment $-d_k$, and $d_k$ equals the DI change less
$\sum_j\Z_k^\top\Delta\E\Z_j\langle\ubar\rangle_j$, which is zero with mass matching; for a
constant field $d_k=0$ (\cref{prop:theta}(a)).  The force is mass-proportional, so
\cref{prop:increment} bounds its effect by $\|\MM^{-1}\W\|_\infty|d_k|/\Z_k^\top\W\Z_k$, a
uniform velocity of the size of the momentum defect over the phase mass.  Exactly, DP
changes the momentum by $\sum_j\Z^\top\Delta\E\Z_j\langle\ubar\rangle_j$, provided the sweeps
conserve the $\W$-momentum (no clamped or snapped fractions).
\end{proof}

\emph{Derivation check.} An independent check of the propositions of
\cref{sec:constant,sec:obstruction,sec:designs} and of \eqref{eq:match} (randomised NumPy
checks, 653 passing) is archived with the proof reviews of the audit, folder
\code{mimetic-momentum-transport}; its corrections
(identified values in \cref{prop:constant}(b), the rank of \eqref{eq:match}, the wording of
\cref{prop:flux} and of the trilemma, the per-component condition for a local $\Y$, the
measured, not analytic, cut-volume independence of the increment bound) are incorporated.

%======================================================================
\section{Results}\label{sec:results}
All runs: uniform periodic meshes, \code{MimeticProjection}, the default Lax--Wendroff
transport with the CTU switch, drivers in \code{mm-runs/mm\_momentum/} (\code{freestream.jl},
\code{smallcell.jl}, \code{drop.jl}, \code{geometry.jl}, \code{smoke.jl}; outputs \code{*.txt}).
CFL is the fraction of the full-cell VOF bound of \code{stable\_timestep} (the WY bound admits
$0<\mathrm{CFL}\le\tfrac12$).

\subsection{Free stream}\label{sec:freestream}
\meas\ Uniform $U\in\{(1,0.37),(0,1)\}$ in both phases, a drop ($R=0.2$, cap $10^{-6}$ of a
cell entering a new row) or a band between the lines $x+2y=c\pm0.17$ (one line $10^{-6}h$
from a grid vertex), $n=32$, 30 steps, $\rho_\gas\in\{1,10^{-3}\}$, CFL $0.1$ and $0.45$, cell
fraction floor $\sqrt\epsilon$ and $10^{-12}$; maximum over the runs:
\begin{center}\footnotesize
\setlength\tabcolsep{4pt}
\begin{tabular}{@{}llrrrr@{}}
\toprule
case & design & $\frac{\max|u-U|}{|U|}$, transport & after the step & $\frac{\max|\Delta P_{\mathrm{mix}}|}{\rho V|U|}$ & min $\min(\alpha,1-\alpha)$\\
\midrule
drop, $\sqrt\epsilon$ & D0, DI & \g 0 & \g 0 & $2.1\times10^{-4}$ & $1.0\times10^{-7}$\\
 & DP+mm & \g $1.4\times10^{-16}$ & \g $1.4\times10^{-16}$ & \g $3.0\times10^{-12}$ & \\
 & DG+mm & $5.5\times10^{-5}$ & $5.5\times10^{-5}$ & $1.0\times10^{-11}$ & \\
 & DS / DF & \rd $4.6\times10^{3}$ / $8.6\times10^{4}$ & \rd & $1.1$ / $2.2\times10^{-3}$ & \\
drop, $10^{-12}$ & D0, DI & \g 0 & \g 0 & $2.1\times10^{-4}$ & $5.3\times10^{-10}$\\
 & DP+mm & \g $1.4\times10^{-16}$ & \g $1.4\times10^{-16}$ & \g $2.4\times10^{-14}$ & \\
 & DG+mm & \rd 310 & \rd 310 & $2.4\times10^{-5}$ & \\
 & DS / DF & \rd $1.6\times10^{7}$ / $1.2\times10^{5}$ & \rd & $1.9$ / $2.2\times10^{-3}$ & \\
band, $\sqrt\epsilon$ & D0, DI & \g 0 & \g 0 & $3.7\times10^{-8}$ & $1.3\times10^{-7}$\\
 & DP+mm & \g $2.6\times10^{-16}$ & \g $2.6\times10^{-16}$ & $2.4\times10^{-8}$ & \\
 & DG+mm & $1.8\times10^{-3}$ & $1.8\times10^{-3}$ & $2.4\times10^{-8}$ & \\
 & DS / DF & \rd $6.0\times10^{3}$ / $4.7\times10^{5}$ & \rd & $0.23$ / $1.4\times10^{-2}$ & \\
\bottomrule
\end{tabular}
\end{center}
D0 and DI keep constants bitwise (the $\W$ sweeps are in fluctuation form and the DI right
side is exactly zero), through every sweep and every step, at all CFL numbers and cut volumes
tried; DP to round-off.  Their $\MM$-momentum of the constant field drifts by the obstruction
of \cref{prop:obstruction} ($2\times10^{-4}$ in 30 steps for the drop) unless the pressure
points are mass-matched, after which it is round-off; the band case keeps $2.4\times10^{-8}$,
\infr\ the VOF's own volume change by fraction snapping (an $O(\text{floor})$ effect, the same
for every design).  DG preserves constants only up to round-off amplified by the inverse
inertia of slaved values (the cancellation in \cref{prop:redistribution} is exact in
arithmetic, not in floating point); DS and DF do not preserve constants
(\cref{prop:flux,prop:theta}).  The CFL $0.9$ runs of \code{freestream.txt} are rejected by the
WY bound check, not by the designs.

\subsection{Small cut volumes and the time step}\label{sec:smallcell}
\meas\ A drop carried by $U=(1,0.37)$ plus a solenoidal vortex of amplitude $0.3$ crosses
the box once ($t=1$, $n=32$), so the interface sweeps through near-empty faces.  Steps to
$t=1$, maximum phase speed over the run relative to the initial maximum on values whose phase
aperture is at least $0.05$ (``wet'', the floor of the static-drop test) and over all values
(``all''); a run stops ($\dagger$) when $E>100E_0$ or a value exceeds $10^4u_0$:
\begin{center}\scriptsize
\setlength\tabcolsep{2.5pt}
\begin{tabular}{@{}lll|rrr|rrr|rrr|rr@{}}
\toprule
 & & & \multicolumn{3}{c|}{D0 (current)} & \multicolumn{3}{c|}{DI} & \multicolumn{3}{c|}{DP + mass matching} & DG+mm & DF\\
$\rho_\gas$ & floor & CFL & steps & wet & all & steps & wet & all & steps & wet & all & \multicolumn{2}{c}{stopped at step}\\
\midrule
1 & $\sqrt\epsilon$ & 0.10 & 538 & 1.01 & 1.01 & 542 & 1.03 & 2.43 & 541 & 1.02 & 1.14 & 3 & 126\\
1 & $\sqrt\epsilon$ & 0.25 & 215 & 1.01 & 1.01 & 216 & 1.02 & 1.05 & 216 & 1.03 & 1.05 & 1 & 24\\
1 & $\sqrt\epsilon$ & 0.50 & 108 & 1.01 & 1.00 & 108 & 1.02 & 1.04 & 108 & 1.02 & 1.05 & 21 & 42\\
1 & $10^{-12}$ & 0.10 & 538 & 1.01 & 1.01 & 542 & 1.03 & 2.43 & 541 & 1.02 & 1.15 & 3 & 54\\
1 & $10^{-12}$ & 0.25 & 215 & 1.01 & 1.01 & 216 & 1.02 & 1.05 & 216 & 1.03 & 1.05 & 1 & 43\\
1 & $10^{-12}$ & 0.50 & 108 & 1.01 & 1.00 & 108 & 1.02 & 1.04 & 108 & 1.02 & 1.05 & 7 & 55\\
$10^{-3}$ & $\sqrt\epsilon$ & 0.10 & 735 & 6.13 & 204 & 1228 & 1.69 & 200 & 1177 & 1.79 & 159 & 5 & 944\\
$10^{-3}$ & $\sqrt\epsilon$ & 0.25 & 271 & 3.74 & $1.0\times10^{3}$ & 708 & 1.76 & 171 & 610 & 2.19 & 162 & 1 & 5\\
$10^{-3}$ & $\sqrt\epsilon$ & 0.50 & 124 & 5.84 & 610 & 363 & 3.35 & 263 & 402 & 2.03 & 110 & 6 & 2\\
$10^{-3}$ & $10^{-12}$ & 0.10 & 827 & 5.30 & 637 & 1440 & 1.88 & 370 & 1500 & 1.71 & $5.2\times10^{3}$ & 5 & 120\\
$10^{-3}$ & $10^{-12}$ & 0.25 & 336 & 6.33 & 360 & 950 & 2.60 & $1.3\times10^{3}$ & 822 & 2.11 & 373 & 1 & 5\\
$10^{-3}$ & $10^{-12}$ & 0.50 & \rd 145$^\dagger$ & \rd 6.13 & \rd $1.0\times10^{4}$ & 491 & 1.81 & $1.9\times10^{3}$ & 622 & 5.31 & 352 & 6 & 2\\
\bottomrule
\end{tabular}
\end{center}

$\MM$-energy maximum over the run relative to its initial value: D0 $1.0001$--$1.008$, DI
$1.0008$--$1.0096$, DP+mm $\le1.0000$ (to the four printed digits in every run).
\begin{itemize}[leftmargin=*,itemsep=1pt]
\item DI and DP run at every CFL up to the WY maximum $\tfrac12$ with cut volumes down to
  $10^{-12}$: no small-cell restriction, as \cref{prop:increment} predicts.  DG and DF fail
  in 1--126 steps at every CFL (\cref{rem:unbounded}); smaller steps do not help, because the
  metric jumps do not scale with $\dt$.
\item At $\rho_\gas=1$ the values stay bounded for all three (wet $\le1.03$); DI adds new
  extrema on small-aperture values (all $\le2.43$), the non-monotone $\MM^{-1}\W$ of
  \cref{prop:increment}.
\item At $\rho_\gas=10^{-3}$ the current method already has gas values of $10^2$--$10^4u_0$ on
  faces with aperture below $0.05$ (the known light-phase velocities on small apertures of
  the projection, \code{report\_mmd.md}); D0 stopped on one at CFL $0.5$, floor $10^{-12}$.
  DI and DP lower the wet maxima (1.7--3.4 and 1.7--5.3 against 3.7--6.3 for D0) but do not
  bound the small-aperture values: DP reached $5.2\times10^3u_0$ at CFL $0.1$, floor $10^{-12}$
  (\textbf{open}).
\item At $\rho_\gas=10^{-3}$ DI and DP take 1.6--3.4 times more steps than D0 for the same
  CFL fraction: the bound of \code{stable\_timestep} is the full-cell VOF bound evaluated with
  the phase velocities, and larger gas values on partially wet faces lower it \infr.  At
  $\rho_\gas=1$ the step counts agree to 1\%.
\end{itemize}

\subsection{Static and translating drop}\label{sec:drop}
\meas\ \code{static\_droplet\_case(:mimetic, ...)} of the test, $t=0.4$; and the same drop in
a periodic box, at rest or translating with $U=(1,0.37)$, $\max|u-U|$ on mixed faces with
phase aperture above $0.05$, first half $\to$ second half of the run ($n=32$):
\begin{center}\footnotesize
\setlength\tabcolsep{4pt}
\begin{tabular}{@{}lllll@{}}
\toprule
 & D0 & D0+mm & DI & DP+mm\\
\midrule
static DR 1, La $\infty$: $E_{\text{late}}/E_{\text{early}}$ & 2.18 & 2.18 & 2.21 & 2.21\\
static DR 1000, La $\infty$: $E_{\text{late}}/E_{\text{early}}$ & 1.68 & 1.64 & 1.71 & 1.67\\
static DR 1, La 1200: gas & $4.11\to7.27\,(10^{-3})$ & $4.11\to7.26$ & $4.11\to7.26$ & $4.11\to7.28$\\
static DR 1000, La 1200: gas & $1.71\to2.86\,(10^{-2})$ & $1.70\to2.86$ & $1.70\to2.88$ & $1.70\to2.89$\\
periodic DR 1, $U=0$ & $7.8\to9.6\,(10^{-3})$ & $7.8\to9.6$ & $7.8\to9.6$ & $7.8\to9.6$\\
periodic DR 1, $U=(1,0.37)$ & $0.191\to0.178$ & $0.191\to0.178$ & $0.179\to0.174$ & $0.179\to0.174$\\
periodic DR 1000, $U=0$ & $1.7\to2.8\,(10^{-2})$ & $1.7\to2.8$ & $1.7\to2.9$ & $1.7\to2.8$\\
periodic DR 1000, $U=(1,0.37)$ & $0.49\to1.78$ & $0.49\to1.78$ & $0.76\to1.32$ & $0.74\to1.56$\\
\bottomrule
\end{tabular}
\end{center}
No design changes the static drop materially (its growth at $n=32$ is the capillary
instability of \code{static\_drop\_stability}, sec.~zero).  Translation raises the spurious
currents 20-fold (DR~1) and 60-fold (DR~1000) for \emph{every} design, the current one
included: the loss of Galilean invariance does not come from the metric of the transport
\infr\ (constants are transported exactly, \cref{sec:freestream}); it is open.

\subsection{Cost}\label{sec:cost}
\meas\ ($n=32$, drop, \code{smoke.jl} and \code{smallcell.jl}): the band has 324--354 values;
one solve with the existing band factorisation takes 6--10\,$\mu$s, the snapshot and the
products of the prototype (which rebuilds the sparse metric each step) 0.6\,ms; the LSQR of
DG 2--3.5\,ms.  Wall time per step with the diagnostics of the drivers: D0 8.0--9.5\,ms, DI
8.6--10.3\,ms, DP 8.8--11.2\,ms.  In production the extra cost is one back-substitution and
two metric products per step (the factorisation is shared with the viscous step and the
projection).

\subsection{Not measured (stopped by the owner)}\label{sec:missing}
The KH (inviscid, momentum and energy drift, energy versus $\dt$), RT DR3/DR1000 $n=32$,
capillary waves DR1/DR1000 $n=16/32/64$, combined case and the per-case timing runs were
prepared (\code{cases.jl}, \code{queue.sh}) but not run: the study was parked.  The energy
statement of R3 is therefore only analysed (\cref{prop:theta,rem:unbounded}), not measured
as $\dt\to0$.

\subsection{Answers}\label{sec:answers}
\begin{enumerate}[leftmargin=*,itemsep=2pt]
\item \emph{R1 and R2 together} are impossible for the Gram metric as assembled
  (\cref{prop:obstruction}): its mass differs from the VOF mass by $O(h^2)$ and moves with
  the interface.  Mass matching of the pressure points ($O(h^3)$ slides, one $3\times3$
  solve) removes the obstruction; the static drop is unchanged by it to 3\%.
\item \emph{Flux forms} with the VOF fluxes cannot preserve constants (\cref{prop:flux}):
  per value the Gram metric and the VOF masses differ by $O(1)$.
\item \emph{Bounded designs.} Any design that adds the metric change $\E^+-\E$ to the momentum
  (DF, DS, DG) is unbounded at slaved values and fails R4 for every $\dt$.  The increment
  form DI is bounded independently of the cut volume ($\lambda_{\max}(\W,\MM)\le2.7$), exact
  for constants in every sweep, and costs one back-substitution; DP adds exact $\MM$-momentum
  with mass matching.
\item \emph{R3} (exact semi-discrete energy) needs $\theta=\tfrac12$ or the skew
  redistribution, i.e.\ explicit $\E$ terms: not reachable with a bounded design for the
  present metric, whose assembly changes discretely in time (status changes of values).
  \textbf{Open}: a metric that is continuous in time at the level of single values (no donor
  switches, continuous slaving) is the prerequisite for DS/DG.
\item \emph{Open:} gas values on faces with aperture below $0.05$ at density ratio 1000 are
  unbounded with D0, DI and DP alike; Galilean invariance of the capillary drop.
\end{enumerate}


%======================================================================
\section{Aperture activity}\label{sec:aperture}
The negative results above come from values that the activity rule of the package creates:
a phase value $(f,\pi)$ is active iff its staggered phase volume is positive
(\code{\_phase\_face\_is\_active}).  A value with volume but without a wetted face
(aperture $a^\pi_f=0$, a \emph{slaved} value) carries no flux, enters $\MM$ only through a tie
to a donor with the weight of its staggered mass (down to $10^{-14}$), and is where DF, DS and
DG blew up (\cref{rem:unbounded}).  This section studies the alternative: \emph{a value is
active iff its phase aperture is positive}, the flux carriers of the projection.

\subsection{Formulation}\label{sec:ap:formulation}
\emph{Unknowns and metric.}  The degrees of freedom are the values with $a^\pi_f>0$ (and the
identified grid-line values).  The metric is the Gram assembly without slaved values: per cut
cell and phase the block on the wetted sides, $\tfrac12\rho^\pi h^2$ per side of an uncut
cell, lumped masses where no block gives inertia, and optionally the ties of the values with
$a<\tfrac12$ (``small ties'', weight $(1-2a)m_f$).  The prototype switch
\code{\_PHASE\_ACTIVITY[] = :aperture} applies the rule consistently to the marking of the
phase velocities (absent values are \code{NaN}), the metric assembly (hence the projection,
the capillary force and the viscous step, which act on its degrees of freedom), the
preparation of the $\W$ transport (aperture-inactive values are extended like absent ones)
and its velocity check; \code{\_MIMETIC\_SMALL\_TIES[]} switches the small ties.  The default
path is unchanged.

\emph{Mass and momentum of a cut cell} live in its Gram blocks: $\Z^\top\MM\Z=\rho\sum_c\G_c$
does not change with the activity rule, because slaved values had $(\MM\Z)_f=0$ (ties vanish on
constants).  With mass matching it is $\rho^\pi V^\pi$ (\cref{sec:constant}).  A region with
phase volume but no wetted face of that phase (a cell with all apertures $0$ or $1$ but
$0<\alpha_c<1$, the \code{crossed==0} case, a grid-line face, a sliver behind a trace) has no
value: its mass enters the totals through mass matching (the uncut cell contributes
$\G_c=h^2\mathsf I$ or $0$, the difference to $\alpha_ch^2\mathsf I$ is absorbed by the slides of
the cut cells), and its momentum is that of the neighbouring values.

\emph{Merge of the staggered masses.}  The transport fluxes are still those of the VOF volumes
mapped to the staggered control volumes, i.e.\ they move the staggered masses $m_f$, also of
values that are no longer unknowns.  The prototype merges them: a value keeps
$k_f=\min(m_f,\kappa\MM_{ff})$, and the excess $m_f-k_f$ (all of $m_f$ for a value without
degree of freedom) goes to the value of the same phase in the face ring (same axis) with the
largest aperture.  This defines $\widetilde\W=\operatorname{diag}(\widetilde m)$ on the degrees of
freedom, with $\Z^\top\widetilde\W\Z=\Z^\top\W\Z$ (donors of the same axis and phase) and
$\widetilde m_f\le\kappa\MM_{ff}$ for every value that is not a donor; the same merge maps the
covector $m_f\Delta u_f$ of the sweeps.  All designs of \cref{sec:designs} are defined with
$\widetilde\W$ in place of $\W$ ($\kappa=1$).

\subsection{Analysis}\label{sec:ap:analysis}
\begin{proposition}[What the activity rule changes and what it does not]\label{prop:aperture}
Under aperture activity:
\begin{enumerate}[leftmargin=*,itemsep=1pt,label=(\alph*)]
\item \emph{Births and deaths (without small ties).}  For aperture data that close a polygon
  (a straight interface in the cell) the row of a value with wetted length $b$ is $O(b)$ and
  its diagonal between $b^2$ and $b/2$ (the clamp margin gives exactly $0.01\,b\rho h^2$ for a
  single partial side): a value is born and dies with a vanishing row, and a birth or death
  changes $\MM$-momentum and $\MM$-energy by $O(b)$ (not $O(\MM_{ff})$: the off-diagonal
  entries are $O(b)$ while $\MM_{ff}$ may be $O(b^2)$).  The slaved values and their ties are
  gone.  The metric is \emph{not} continuous in time, however: the block of the other phase
  jumps when the cell becomes uncut, the pressure point jumps for non-closed (averaged-trace)
  data, whose rows need not be $O(b)$, and $\widetilde\W$ jumps at the lump threshold; donor
  switches of the merge remain.  With the small ties a newborn value enters with a tie of
  weight $(1-2a)m_f\to m_f$.
\item \emph{Inertia.}  $\lambda_{\max}(\W,\MM)\ge m_f(\MM^{-1})_{ff}\ge m_f/\MM_{ff}$, which is
  unbounded at birth without small ties ($\MM_{ff}=O(b^2)$ for a corner wedge while $m_f$ need
  not vanish).  The merge does not cure it: the cap $\widetilde m_f\le\kappa\MM_{ff}$ bounds the
  diagonal ratio, not $\lambda_{\max}(\widetilde\W,\MM)$, which grows like $10^{-3}/b$ because a
  corner wedge makes $\MM$ nearly singular ($\lambda_{\min}\approx3.5\,b^2$).  The small ties
  restore an inertia of order $m_f$ for $a<\tfrac12$ (measured $\lambda_{\max}\le2.5$ with them).
\item \emph{The algebra.}  \Cref{prop:obstruction} carries over (the merge keeps the totals,
  the VOF mass is conserved), and so does the conclusion of \cref{prop:flux}.  \Cref{prop:theta}(a)
  (constants) holds with $\widetilde\W$; (b) and (c) acquire the merge term
  $\mu(v)=\sum_f(m_f-k_f)(v_{d(f)}-v_f)\Z_f$ (the excess of $f$ moves with the velocity of its donor
  in $\widetilde\W u$ but with its own in the merged flux covector), so momentum and energy are
  no longer exact for DF, DS and DG; the merged covector cannot satisfy both
  $\Z^\top\widetilde\C=0$ and $\dt\widetilde\C(\Z)=-\Delta\widetilde\W\Z$ because the keep/excess split
  changes between the levels.  The momentum restoration of DP (and DGP, DG with the same
  restoration) uses the momentum $\Z^\top(\MM^+\ubar+\text{rhs})$ of the actual update and is
  therefore exact whatever the update, $\mu$ included, up to the constant part
  $\Z^\top\Delta\MM\Z\langle\ubar\rangle$ (zero with mass matching).
\end{enumerate}
\end{proposition}
\begin{proof}
(a) For closed data the block is that of MM on a polygon; its rows are $b_k(x_k-x^*)$ and the
stabilisation, both $O(b)$; the counterexamples (non-closed data with an entry $2.07$ at
$b=10^{-8}$; the cut$\to$uncut jump $1.48$ of the other block; the lump threshold) are in the
derivation check.  The momentum change of a dying value is $\sum_g\MM_{fg}u_g\,\Z_g$-weighted,
$O(b)$.  (b) The diagonal of $\W^{1/2}\MM^{-1}\W^{1/2}$.  (c) Expand the update with
$\widetilde\W$ and compare the merged covector $\sum_f k_f\Delta u_f+\sum_f(m_f-k_f)\Delta u_f$
(at the donors) with $\widetilde\W\Delta u$.  DP: \cref{sec:designs}, with
$\MM^+u^+=\MM^+\ubar+\text{rhs}$.
\end{proof}
The derivation check of this section (379 randomised checks, Python port of the Gram block and a
lattice model) is archived with the proof reviews of the audit (folder
\code{mimetic-momentum-}\allowbreak\code{transport-aperture}); its corrections are incorporated above.  It also
found that an identified grid-line donor moves liquid excess into the gas column, and that a
value without donor keeps $\widetilde m_f>\kappa\MM_{ff}$ (or drops its mass if it has no
degree of freedom); both are conditions of the merge, not measured to matter.
So aperture activity can only help where the old failures were caused by slaved values; it
does not by itself make the metric continuous in time.

\subsection{Measurements under aperture activity}\label{sec:ap:results}
Same drivers as \cref{sec:results} with \code{MM\_ACTIVITY=aperture} (outputs
\code{*\_ap.txt}; \code{\_apnt}: without small ties).

\emph{Free stream} \meas\ (\code{freestream\_ap.txt}, same cases as \cref{sec:freestream}):
\begin{center}\footnotesize
\setlength\tabcolsep{4pt}
\begin{tabular}{@{}lrrrr@{}}
\toprule
design & drop, $\sqrt\epsilon$ & drop, $10^{-12}$ & band & $\max\frac{|\Delta P_{\mathrm{mix}}|}{\rho V|U|}$ (drop)\\
\midrule
D0, DI & \g $5.2\times10^{-17}$ & \g $5.2\times10^{-17}$ & \g 0 & $2.2\times10^{-4}$\\
DP+mm & \g $5.2\times10^{-17}$ & \g $5.2\times10^{-17}$ & \g $2.6\times10^{-16}$ & \g $3.1\times10^{-12}$\\
DG+mm & \g $4.1\times10^{-14}$ & \g $1.4\times10^{-14}$ & \g $2.0\times10^{-15}$ & \g $3.1\times10^{-12}$\\
DS & \rd $5.1\times10^{4}$ & \rd $2.6\times10^{5}$ & \rd 140 & $2.0\times10^{-2}$\\
DF & \rd $3.3\times10^{4}$ & \rd $4.3\times10^{6}$ & \rd 460 & $2.2\times10^{-2}$\\
\bottomrule
\end{tabular}
\end{center}
(maximum $|u-U|/|U|$ after the step).  DG with mass matching is now exact for constants to
round-off ($\le4\times10^{-14}$, against up to 310 under volume activity) and conserves the
$\MM$-momentum of the constant field: the $\theta=\tfrac12$ defect $\tfrac12\Delta\E\Z$ is
redistributed exactly, and no slaved inertia amplifies the round-off.  DS and DF still do not
preserve constants (\cref{prop:aperture}(c)).

\emph{One-step identities} \meas\ (\code{smoke\_ap.jl}, drop, $\rho_\gas=10^{-3}$, eight steps,
uniform field and the vortex field of \cref{sec:smallcell}): DG+mm: constants $6\times10^{-16}$,
momentum change of the constant field $1.9\times10^{-16}$, energy change $-1.4\times10^{-17}$;
with the vortex the largest correction $|u^+-u_\W|$ is $0.16$ (volume activity: $4\times10^3$ to
$6\times10^4$).  DI and DP: $0.03$.  DS, DF: $23$ and $18$, on small-aperture values.

\emph{Small cut volumes} \meas\ (\code{smallcell\_ap.txt}, \code{smallcell\_apnt.txt}; runs
completed out of six per density, maxima over the completed ones):
\begin{center}\footnotesize
\setlength\tabcolsep{4pt}
\begin{tabular}{@{}l|cc|cc|cc@{}}
\toprule
 & \multicolumn{2}{c|}{volume activity} & \multicolumn{2}{c|}{aperture, small ties} & \multicolumn{2}{c}{aperture, no small ties}\\
design & $\rho_\gas=1$ & $10^{-3}$ & $\rho_\gas=1$ & $10^{-3}$ & $\rho_\gas=1$ & $10^{-3}$\\
\midrule
D0 & \g 6/6, 1.01 & 5/6, 6.33 & \g 6/6, 1.03 & \g 6/6, 3.35 & \rd 0/6 & \rd 1/6\\
DI & \g 6/6, 1.03 & \g 6/6, 3.35 & \g 6/6, 1.03 & \g 6/6, 2.28 & \rd 0/6 & \rd 0/6\\
DP+mm & \g 6/6, 1.03 & \g 6/6, 5.31 & \g 6/6, 1.03 & \g 6/6, 1.89 & \rd 0/6 & \rd 0/6\\
DG+mm & \rd 0/6 (1--21) & \rd 0/6 (1--6) & \rd 1/6 (85--452) & \rd 0/6 (46--761) & \rd 0/6 & \rd 0/6\\
DS & --- & --- & \rd 0/6 & \rd 0/6 & \rd 0/6 & \rd 0/6\\
DF & \rd 0/6 & \rd 0/6 & \rd 0/6 & \rd 0/6 & \rd 0/6 & \rd 0/6\\
\bottomrule
\end{tabular}
\end{center}
Entries: completed runs, worst ``wet'' maximum $\max|u|/u_0$; in parentheses the steps at
which DG stopped.  The ``all'' maxima (small-aperture values included) of the completed runs at
$\rho_\gas=10^{-3}$: D0 $897$, DI $1.0\times10^3$, DP $2.6\times10^3$ (volume activity: $\ge10^3$,
$1.9\times10^3$, $5.2\times10^3$); $\MM$-energy maxima $\le1.0018E_0$ for D0, DI, DP.
\begin{itemize}[leftmargin=*,itemsep=1pt]
\item With the small ties, aperture activity keeps D0, DI and DP bounded at every CFL up to
  $\tfrac12$ and cut volumes to $10^{-12}$, and lowers their wet maxima at $\rho_\gas=10^{-3}$
  (D0 $3.35$ against $6.33$; DP $1.89$ against $5.31$).
\item DG survives an order of magnitude longer than under volume activity but is still
  unbounded; DS and DF are unbounded.  The remaining jumps of the metric (the small ties with
  their donors, the merge donors, lumping) are enough.
\item Without the small ties \emph{every} design fails, the current transport included: values
  born with Gram inertia $O(b)$ receive pressure and force increments of order $1/b$ from the
  projection (\cref{prop:aperture}(a,b)).  The continuity of (a) is therefore not usable with the
  present Gram blocks: the small ties are needed for the projection itself, and they bring the
  birth jumps back.
\end{itemize}

\emph{Static and translating drop} \meas\ (\code{drop\_*\_ap.txt}, $n=32$, first half $\to$ second
half): D0, DI and DP+mm under aperture activity keep the inviscid static energy growth
($2.18$--$2.22$ at DR~1, $1.69$--$1.72$ at DR~1000) and the DR~1 currents (viscous gas
$4.1\to6.1\times10^{-3}$ against $4.1\to7.3\times10^{-3}$).  At DR~1000 the viscous static gas
current doubles ($3.9\to5.8\times10^{-2}$ against $1.7\to2.9\times10^{-2}$; Ca $7.4\times10^{-4}$,
below the test bound $5\times10^{-3}$) and the translating drop is 20--60\% worse (D0
$0.60\to2.32$ against $0.49\to1.78$).  DG+mm completes the static drops (inviscid DR~1 growth
$2.42$) and the DR~1 translating drops but stops with a non-finite velocity in the DR~1000
translating drop at $n=16$.


\subsection{\texorpdfstring{The aperture definition: plain $\tfrac12$--$\tfrac12$ average}{The aperture definition: plain half-half average}}\label{sec:ap:plain}
\emph{How the apertures are defined now.}  The metric, the activity rule, the projection and the
VOF sweeps (which freeze a copy at the start of the step) all use the interface-owned phase
apertures of \code{\_current\_phase\_face\_apertures!}, i.e.\ \code{\_phase\_face\_apertures}: the
liquid aperture is the mean of the two cells' traces, $a_f=\tfrac12(t_{c_-}+t_{c_+})$, with the
trace $t_c$ of \code{\_face\_aperture} the part of the face below the PLIC line of $c$ ($0$ for
$\alpha_c\le0$, the full open face for $\alpha_c\ge1$, a solid clipped first), and the gas
aperture $1-a_f$.  The exceptions that clip to zero or one are:
\begin{enumerate}[leftmargin=*,itemsep=1pt,label=(\roman*)]
\item a phase absent in both cells gets aperture exactly $0$ (the other phase the open face);
  this coincides with the plain mean, since both traces are then $0$ or full;
\item the round-off snap \code{\_snap\_phase\_apertures}: an aperture below $64\epsilon$ is set to
  $0$ (and the other to the open face), against $10^{-17}$ openings left by two independently
  clipped polygons;
\item indirectly, the VOF floor: cell fractions within \code{fraction\_floor} ($\sqrt\epsilon$) of
  $0$ or $1$ are snapped after each sweep, so their traces become $0$ or full;
\item in the metric only, a grid-line face between a pure liquid and a pure gas cell is
  identified (one value, each cell sees its own side fully wet), whatever its mean aperture
  ($\tfrac12$);
\item the boundary faces (\code{\_phase\_boundary\_face\_apertures}) use the single trace with
  the same snap.
\end{enumerate}
No pure-neighbour rule is left in this branch (it was removed with MMd, \code{report\_mmd.md}).
The present aperture \emph{is} therefore the plain $\tfrac12$--$\tfrac12$ mean up to (ii)
(round-off) and the grid-line identification (iv).  The prototype switch
\code{\_APERTURE\_PLAIN[]} (environment \code{MM\_APERTURE=plain}) removes (i) and (ii); (iii)
belongs to the VOF and (iv) to the metric and are kept.

\emph{What changes in principle.}  A value is born when either neighbour's trace reaches the face
and dies when neither does --- the same events as now, since (i) is not a clip.  Without (ii) a
phase can be active on a $10^{-17}$ opening: such values have an inertia of order $b\le10^{-16}$
and, under aperture activity, become degrees of freedom; with the small ties their inertia is
that of their staggered mass instead.  The traces move continuously while one PLIC line
translates, but jump when the reconstruction changes its normal between steps (a new Youngs or
LHF normal pivots the line), so births and deaths are not smoother in time with either
definition: the time continuity is set by the reconstruction, not by the averaging.  The Gram
block of a cell uses the shared averaged apertures of its four sides, which need not close a
polygon with any one cell's PLIC (\code{metric\_options}); the VOF volume fluxes are geometric
(donating regions of the PLIC of the upwind cell), not apertures, so neither definition makes the
fluxes and the block's wetted sides consistent.  The identities of \cref{sec:ap:analysis} do not
depend on how the apertures are computed.

\emph{Time continuity} \meas\ (\code{transitions.jl}, aperture activity, $\rho=(1,10^{-3})$,
$U=(1,0.37)$, $n=32$, 60 steps at CFL $0.45$; per step):
\begin{center}\footnotesize
\setlength\tabcolsep{4pt}
\begin{tabular}{@{}llrrrrrr@{}}
\toprule
case & aperture & births & deaths & $a$ at birth (median, min) & $a$ at death (median, min) & $\max|\Delta a|$ & tie switches\\
\midrule
drop & current & 17.9 & 18.0 & 0.20, $2.2\times10^{-4}$ & 0.19, $1.9\times10^{-4}$ & 0.997 & 29.9\\
drop & plain & 17.9 & 18.0 & 0.20, $2.2\times10^{-4}$ & 0.19, $1.9\times10^{-4}$ & 0.997 & 29.9\\
band & current & 73.6 & 72.5 & 0.20, $3.7\times10^{-4}$ & 0.17, $8.6\times10^{-3}$ & 0.572 & 136.0\\
band & plain & 73.1 & 73.1 & 0.20, $3.7\times10^{-4}$ & 0.17, $1.1\times10^{-16}$ & 0.572 & 136.5\\
\bottomrule
\end{tabular}
\end{center}
($\max|\Delta a|$: largest change in one step of a value that stays active; no value had
$0<a<10^{-6}$.)  The births and deaths of aperture-active values coincide with the births and
deaths of staggered volume (drop: $18.0$ per step): a phase reaches a face exactly when it enters
an adjacent cell, and the PLIC of that cell then already has a trace of typically $0.2$ on the
face.  A value is therefore born with a \emph{finite} aperture, not with $a\to0$, and the
apertures of surviving values jump by up to $1$ in one step (the reconstructed line of a cell
changes its normal from step to step).  The plain average changes nothing measurable: identical
counts for the drop, one death on a $10^{-16}$ opening in the band.  The time continuity
required by \cref{prop:aperture}(a) is not a property of the PLIC apertures with either
definition.

\emph{Key checks with the plain average} \meas: free stream (\code{freestream\_ap\_pl.txt}),
small-cell sweep (\code{smallcell\_ap\_pl.txt}) and the one-step identities
(\code{smoke\_ap\_pl.txt}) are identical to those of the current aperture to the printed digits
(the small-cell step counts agree exactly), except that DP's restoration now acts on round-off
($\le2.4\times10^{-14}$ for constants, momentum change $1.9\times10^{-14}$ with the vortex).
The impact with D0 is bitwise the same run (3092 steps, same maxima).


\subsection{Answers for aperture activity}\label{sec:ap:answers}
\emph{Drop impact} \meas\ (\code{impact.jl}: pool $y<0.3$, drop $R=0.08$ falling at $2$, DR~1000,
$\sigma=0.05$, $\mu=(10^{-3},10^{-6})$, $n=32$, $t\le0.25$, package CGFM force; $\max|u|$ over all
phase values with its location, $\MM$-energy relative to the start):
\begin{center}\footnotesize
\setlength\tabcolsep{4pt}
\begin{tabular}{@{}lrrll@{}}
\toprule
run & steps & $\max E/E_0$ & $\max|u|$ & where\\
\midrule
D0, volume activity & 779 & \rd 186 & $4.6\times10^4$ & slaved gas value ($a=0$, $\alpha=1.1\times10^{-8}$)\\
D0, aperture & 3092 & 1.12 & $1.2\times10^4$ & gas, $a=5\times10^{-5}$, $\alpha=7\times10^{-9}$\\
DI, aperture & --- & \rd $10^{14}$ & $5\times10^{8}$ & blew up at $t=0.207$\\
DP+mm, aperture & 586 & 1.19 & $1.2\times10^4$ & gas, $a=1.3\times10^{-4}$, $\alpha=4\times10^{-8}$\\
DG+mm, aperture & --- & \rd $\infty$ & --- & non-finite at $t=0.065$\\
\bottomrule
\end{tabular}
\end{center}
The plain aperture gives the same runs (D0 and DP bitwise the same maxima and steps, DG non-finite
at the same time).  Aperture activity removes the energy excursion of the current method
($186E_0$, carried by slaved gas values) and the slaved outliers, but the largest gas velocities
($10^4$, against the physical scale $2$) remain, now on values of aperture $10^{-5}$--$10^{-4}$ and
gas fraction $10^{-8}$: the small-aperture half of the outliers found by the separate impact
diagnosis (branch \code{study/impact-velocities}).  DP completes the impact with a fifth of the steps of D0.

\emph{Energy as $\dt\to0$.}  The test of \cref{sec:missing} was attempted on the vortex drop with
Lax--Wendroff \emph{without} the CTU switch (the switch is upwind in the band, an $O(1)$
dissipation): every design, the current one included, left the advective stability bound within
2--60 steps at $\rho_\gas=10^{-3}$ and in most runs at $\rho_\gas=1$ (\code{invariants*.txt}).
The completed runs with aperture activity at $\rho_\gas=1$ give $\sum\Delta E_{\text{transport}}/E_0$
$=5.5\times10^{-5},1.8\times10^{-4},1.9\times10^{-4}$ for DI at $\dt$, $\dt/2$, $\dt/4$ (not
decreasing, as \cref{prop:theta}(c) predicts for $\theta=1$), and $2.2\times10^{-3}$,
$3.2\times10^{-3}$ for DGP.  The same runs confirm exact $\MM$-momentum of DP and DGP with mass
matching under the merge: $|\sum\Delta P|/S\le6\times10^{-11}$ against $10^{-4}$--$10^{-3}$ for D0
and DI.

\emph{Answers.}
\begin{enumerate}[leftmargin=*,itemsep=2pt]
\item \emph{Formulation.}  Aperture activity is implemented consistently (marking, metric,
  projection, capillary force, viscous step, transport preparation, VOF velocity check) as a
  non-default switch; slaved values disappear; regions with volume but no wetted face enter
  through mass matching, their staggered masses through a merge into same-axis donors.  The
  $\MM$-mass of a constant field is unchanged and equals the VOF mass with mass matching.
\item \emph{Births and deaths} happen when a phase enters or leaves an adjacent cell, with a
  finite aperture (median $0.2$) and with aperture jumps of up to $1$ per step for surviving
  values; the metric is not continuous in time (cut$\to$uncut, non-closed data, lumping, donors).
  The small ties are required by the projection itself (without them every design blows up),
  and they bring an $O(m_f)$ jump at every birth.
\item \emph{Identities.}  Exact free stream: D0, DI, DP, DG (to round-off, DG now without
  amplification).  Exact $\MM$-momentum: DP and DGP with mass matching (also under the merge).
  Semi-discrete energy: no bounded design (DG's exactness is lost to the merge term $\mu$, and DG
  is unbounded).  The trilemma of \cref{prop:theta} and the flux-form impossibility are unchanged.
\item \emph{Boundedness and time step.}  With the small ties D0, DI and DP are bounded at every
  CFL up to $\tfrac12$ for cut volumes to $10^{-12}$; DG, DS and DF are not.  DI blew up in the
  impact; DP and D0 completed it.
\item \emph{Plain $\tfrac12$--$\tfrac12$ aperture.}  The package already uses it up to a round-off
  snap and the grid-line identification; removing the snap changes no result.  The time
  continuity is set by the PLIC reconstruction, not by the averaging.
\item \emph{Recommendation} (for the owner): adopt aperture activity only together with the
  small ties, and only as a fix of the current method's slaved-value outliers (energy excursion
  of the impact $186E_0\to1.1E_0$); it does not open the way to an $\MM$-native transport with
  explicit metric-change terms.  If a metric-consistent transport is wanted, DP with mass
  matching remains the candidate (exact constants and $\MM$-momentum, bounded in the sweeps,
  completed the impact); DI alone is not robust in the impact.  A time-continuous metric would
  need apertures (or blocks) that change continuously as the interface moves --- e.g.\ traces
  from a reconstruction whose normal is advanced in time instead of recomputed, or blocks built
  from the VOF donating regions --- which is a geometry question, not a transport one.
\end{enumerate}


\end{document}
